\documentclass[12pt, letterpaper]{article}

\input{template/preamble.tex}

\input{template/listings_config.tex}

% \input{template/header.tex}

\begin{document}

\begin{titlepage}

	\begin{tikzpicture}[remember picture,overlay]
		\node[anchor=north west, inner sep=30] at (current page.north west) {
			
\includegraphics[width=3.5cm, height=4cm]{unam.png}
		};
		% \node[anchor=north west, inner sep=30] at (current page.north west) {
		% 	\includegraphics[width=4cm]{escudoUNAM_negro.png}
		% };

		\node[anchor=north east, inner sep=30] at (current page.north east) {
			
\includegraphics[width=3.5cm, height=4cm]{escudoFI_color.png}
		};
		% \node[anchor=north east, inner sep=30] at (current page.north east) {
		% 	\includegraphics[width=4cm]{escudoFI_negro.png}
		% };
	\end{tikzpicture}

	\begin{center}
		\vspace*{3.5cm}

		\Large
		\textbf{UNIVERSIDAD NACIONAL AUTÓNOMA DE MÉXICO}\\
		\textbf{FACULTAD DE INGENIERÍA}\\
		\textbf{INGENIERÍA EN COMPUTACIÓN}\\

		\vspace*{1.5cm}
		\large
		\textbf{Laboratorio de Organización y Arquitectura de Computadoras}\\
		\large
		\vspace*{1cm}
		\textbf{Reporte - Diseño de Máquinas de Estado}\\
		\vspace{1.5cm}
	\end{center}

	%grupo de la materia
	\large
	\textbf{INTEGRANTES EQUIPO 8:}
	\begin{itemize}
		\item Saldaña Navarrete Andrea - \textit{Grupo Teoria:} 1
		\item Tepal Briseño Hansel Yael - \textit{Grupo Teoria:} 1
		\item Ugartechea González Luis Antonio - \textit{Grupo Teoria:} 1
	\end{itemize}
	\vspace{2.5cm}

	\large
	%fecha de entrega
	\textbf{NOMBRE DELPROFESOR:} Ing. Diana Karen Pérez Pérez \\
	\textbf{GRUPO:} 06\\
	\textbf{SEMESTRE:} \underline{2027-1}\\
	\vspace{0.5cm}

	\begin{center}
		\textbf{FECHA DE ENTREGA:} 10 de septiembre del 2026\\
	\end{center}
\end{titlepage}


\section*{Objetivo}

Analizar, diseñar e implementar una máquina de estados finitos a partir de una carta ASM, identificando
el funcionamiento de los flip-flops tipo D como elementos de memoria. Obtener las funciones de estado
siguiente y salida mediante el llenado de memoria y mapas de Karnaugh, para posteriormente comprobar
su funcionamiento utilizando la FPGA DE10-Lite.

\section*{Introducción}

\section*{Desarrollo}

\subsection*{Parte 1. Análisis y diseño de la máquina de estados}

A partir de la carta ASM proporcionada en la Figura 1:

\begin{enumerate}
    \item Identifique los estados, entradas y salidas de la máquina.
    \begin{itemize}
        \item \textcolor{red}{\textbf{Estados:} El sistema cuenta con 4 estados definidos en los rectángulos: \texttt{EST0}, \texttt{EST1}, \texttt{EST2} y \texttt{EST3}.}
        \item \textcolor{red}{\textbf{Entradas:} Las condiciones de transición (rombos) evalúan dos variables: $X$ e $Y$.}
        \item \textcolor{red}{\textbf{Salidas:} Dentro de los bloques de estado se activan las señales $S_1$ y $S_0$. Al estar dentro de los rectángulos, se trata de una máquina de Moore.}
        \begin{itemize}
            \item \textcolor{red}{En EST0: Se activa $S_1$.}
            \item \textcolor{red}{En EST1: Se activa $S_0$.}
            \item \textcolor{red}{En EST2: No hay salidas activas.}
            \item \textcolor{red}{En EST3: Se activan $S_1$ y $S_0$.}
        \end{itemize}
    \end{itemize}
        \item Determine el número de bits necesarios para representar los estados y proponga una codificación.
    \par\textcolor{red}{Dado que se tienen 4 estados, se calcula el número de bits $n$ necesarios usando $2^n \geq 4$, resultando en $n = 2$ bits. Se utilizarán dos variables de estado ($Q_1$ y $Q_0$) que corresponden a dos Flip-Flops tipo D. La codificación binaria propuesta es:}
    \begin{itemize}
        \item \textcolor{red}{\texttt{EST0} = $00$}
        \item \textcolor{red}{\texttt{EST1} = $01$}
        \item \textcolor{red}{\texttt{EST2} = $10$}
        \item \textcolor{red}{\texttt{EST3} = $11$}
    \end{itemize}

    \item Obtenga la tabla de transiciones, considerando el estado actual y las entradas del sistema.
    \par\textcolor{red}{A partir de la carta ASM, la tabla de transiciones muestra el estado siguiente ($Q_1(t+1) Q_0(t+1)$) en función del estado actual ($E_1(t) E_0(t)$) y las combinaciones de las entradas $X$ e $Y$:}
    \begin{center}
        \color{red}
        \begin{tabular}{|c|c|c|c|c|}
        \hline
        \textbf{Estado Actual} & \multicolumn{4}{c|}{\textbf{Entradas $X Y$}} \\
        \cline{2-5}
        $E_1(t) E_0(t)$ & $00$ & $01$ & $11$ & $10$ \\
        \hline
        $00$ (\texttt{EST0}) & $11$ & $11$ & $01$ & $00$ \\
        $01$ (\texttt{EST1}) & $01$ & $11$ & $10$ & $01$ \\
        $10$ (\texttt{EST2}) & $11$ & $11$ & $11$ & $11$ \\
        $11$ (\texttt{EST3}) & $00$ & $00$ & $00$ & $00$ \\
        \hline
        \end{tabular}
    \end{center}

    \item Determine para cada combinación: Estado actual, Entradas, Estado siguiente, Salidas.
    \par\textcolor{red}{Los incisos 3 y 4 se resuelven de manera conjunta en la siguiente tabla de verdad general de la máquina de estados:}
    \begin{center}
        \color{red}
        \begin{tabular}{|cc|cc|cc|cc|}
        \hline
        \multicolumn{2}{|c|}{\textbf{Estado Actual}} & \multicolumn{2}{c|}{\textbf{Entradas}} & \multicolumn{2}{c|}{\textbf{Estado Siguiente}} & \multicolumn{2}{c|}{\textbf{Salidas}} \\
        $E_1(t)$ & $E_0(t)$ & $X$ & $Y$ & $Q_1(t+1)$ & $Q_0(t+1)$ & $S_1$ & $S_0$ \\
        \hline
        0 & 0 & 0 & 0 & 0 & 0 & 1 & 0 \\
        0 & 0 & 0 & 1 & 0 & 0 & 1 & 0 \\
        0 & 0 & 1 & 0 & 0 & 0 & 1 & 0 \\
        0 & 0 & 1 & 1 & 0 & 1 & 1 & 0 \\
        \hline
        0 & 1 & 0 & 0 & 0 & 1 & 0 & 1 \\
        0 & 1 & 0 & 1 & 0 & 1 & 0 & 1 \\
        0 & 1 & 1 & 0 & 0 & 1 & 0 & 1 \\
        0 & 1 & 1 & 1 & 1 & 0 & 0 & 1 \\
        \hline
        1 & 0 & 0 & 0 & 1 & 1 & 0 & 0 \\
        1 & 0 & 0 & 1 & 1 & 1 & 0 & 0 \\
        1 & 0 & 1 & 0 & 1 & 1 & 0 & 0 \\
        1 & 0 & 1 & 1 & 1 & 1 & 0 & 0 \\
        \hline
        1 & 1 & 0 & 0 & 0 & 0 & 1 & 1 \\
        1 & 1 & 0 & 1 & 0 & 0 & 1 & 1 \\
        1 & 1 & 1 & 0 & 0 & 0 & 1 & 1 \\
        1 & 1 & 1 & 1 & 0 & 0 & 1 & 1 \\
        \hline
        \end{tabular}
    \end{center}

    \item Obtenga los mapas de Karnaugh correspondientes a cada entrada de los flip-flops tipo D y a las señales de salida.
    \begin{center}
        \color{red}
        \textbf{Mapa para $Q_1(t+1)$:}\\
        \begin{tabular}{|c|c|c|c|c|}
        \hline
        $E_1 E_0 \backslash X Y$ & 00 & 01 & 11 & 10 \\
        \hline
        00 & 0 & 0 & 1 & 0 \\
        01 & 1 & 1 & 0 & 1 \\
        11 & 0 & 0 & 0 & 0 \\
        10 & 1 & 1 & 1 & 1 \\
        \hline
        \end{tabular}
        \vspace{0.5cm}

        \textbf{Mapa para $Q_0(t+1)$:}\\
        \begin{tabular}{|c|c|c|c|c|}
        \hline
        $E_1 E_0 \backslash X Y$ & 00 & 01 & 11 & 10 \\
        \hline
        00 & 0 & 0 & 0 & 0 \\
        01 & 0 & 0 & 1 & 0 \\
        11 & 0 & 0 & 0 & 0 \\
        10 & 1 & 1 & 1 & 1 \\
        \hline
        \end{tabular}
    \end{center}

    \begin{center}
        \color{red}
        \textbf{Mapa para $S_1$:}\\
        \begin{tabular}{|c|c|c|c|c|}
        \hline
        $E_1 E_0 \backslash X Y$ & 00 & 01 & 11 & 10 \\
        \hline
        00 & 1 & 1 & 1 & 1 \\
        01 & 0 & 0 & 0 & 0 \\
        11 & 1 & 1 & 1 & 1 \\
        10 & 0 & 0 & 0 & 0 \\
        \hline
        \end{tabular}
        \vspace{0.5cm}

        \textbf{Mapa para $S_0$:}\\
        \begin{tabular}{|c|c|c|c|c|}
        \hline
        $E_1 E_0 \backslash X Y$ & 00 & 01 & 11 & 10 \\
        \hline
        00 & 0 & 0 & 0 & 0 \\
        01 & 1 & 1 & 1 & 1 \\
        11 & 1 & 1 & 1 & 1 \\
        10 & 0 & 0 & 0 & 0 \\
        \hline
        \end{tabular}
    \end{center}

    \item Simplifique las funciones booleanas.
    \par\textcolor{red}{De los mapas de Karnaugh anteriores se deducen las siguientes expresiones booleanas simplificadas (suma de productos):}
    \color{red}
    \begin{align*}
        Q_1(t+1) &= E_1 \overline{E_0} + \overline{E_1} E_0 \overline{X} + \overline{E_1} E_0 \overline{Y} + \overline{E_0} X Y \\
        Q_0(t+1) &= E_1 \overline{E_0} + \overline{E_1} E_0 X Y \\
        S_1 &= \overline{E_1} \overline{E_0} + E_1 E_0\\
        S_0 &= E_0\\
    \end{align*}
    \color{black}

    % \item Dibuje el circuito resultante utilizando: Flip-flop tipo D, compuertas y lógica.
    % \par\textcolor{red}{El circuito lógico se implementa conectando las ecuaciones resultantes $D_1$ y $D_0$ a las terminales $D$ de dos flip-flops, empleando compuertas NOT, AND y OR para formar los minitérminos. La salida $S_1$ requiere una compuerta XNOR entre $Q_1$ y $Q_0$, mientras que $S_0$ se toma directamente de la salida del segundo flip-flop. (Se recomienda anexar el diagrama esquemático generado en software en este apartado).}
\end{enumerate}

\begin{figure}[H]
    \centering
    \includegraphics[width=0.4\textwidth]{carta_asm.png}
    \caption{Carta ASM.}
    \label{fig:ASM}
\end{figure}

\section*{Análisis de Resultados}

\begin{figure}[H]
    \centering
    \includegraphics[width=0.5\textwidth]{imagen_ejemplo_1.png}
    \caption{Ejemplo de Simulación.}
    \label{fig:ASM}
\end{figure}

Se logra apreciar que en la simulación con las pruebas correspondientes del estado ...

\section*{Conclusiones}

% Conclusiones por integrante

\input{template/biblio.tex}

\end{document}
